% Table II — the seven-phase pipeline. Used by sections/04-methodology.
% Rule: every phase must declare the evidence it produces AND the condition
% under which it terminates. A phase without a stopping criterion is a task,
% not a method step, and reviewers said so about the previous version.
\begin{table}[!t]
\caption{Reverse-engineering pipeline}
\label{tab:phases}
\centering
\scriptsize
\setlength{\tabcolsep}{3pt}
\begin{tabular}{@{}p{0.15cm}p{1.25cm}p{2.35cm}p{0.5cm}p{2.15cm}@{}}
\toprule
\textbf{\#} & \textbf{Phase} & \textbf{Evidence produced} & \textbf{Hyp.} & \textbf{Stopping criterion} \\
\midrule
1 & Network mapping & Segment map, open ports & --- & All hosts and reachable ports enumerated \\
2 & Introspection & Topic and service inventory & H1 & Inventory stable across sessions \\
3 & Passive broker & Payload classes on the broker & H2 & No command payload after repeated sessions \\
4 & Application audit & Names, types, implicit setup steps & H1 & Every observed name located in code \\
5 & Frame capture & Message framing and ordering & H1 & Client reproduces vendor framing \\
6 & Edge integration & On-device services and bridges & --- & Capability reachable without a host PC \\
7 & Field validation & State transitions, pose deltas & H3, H4 & Effect confirmed on independent topic \\
\bottomrule
\end{tabular}
\end{table}
